\section{Elektronen im Festkörper}
\subsection{Elektronengas}
\begin{itemize}
    \item Valenzelektronen in Metallen quasi frei beweglich
    \item Bewegung der Elektronen in einem Kasten mit unendlich hohen Potentialwänden (Metalloberfläche ist Begrenzung)
    \item Für großes $L$ liegen die Energieniveaus verschiedener $n$ sehr nahe $\Rightarrow$ Leitungsbänder ähnlicher Energienivenaus, dazwischen verbotene Zonen
    \item \begin{equation*}
              E_{n_x,n_y,n_z}=\frac{\hbar^2\pi^2}{2mL^2}(n_x^2+n_y^2+n_z^2)
          \end{equation*}
\end{itemize}

\subsection{Zustandsdichte}
Beschreibt, wie dicht erlaubte Zustände (Leitungsbänder) aneinander gepackt werden können.
\begin{equation*}
    n = \int_0^\infty\underbrace{\left(\frac{2m}{\hbar^2}\right)^{\frac{3}{2}}\frac{\sqrt{E}}{2\pi^2}}_{\text{energet. Zustandsdichte pro Volumen } D(E)}f(T, E)\diff E
\end{equation*}

\begin{sectionbox}
    \subsection{Fermiverteilung $f(T, E)$}
    Besetzungswahrscheinlichkeit für Elektronen im Gleichgewicht. Nähert sich für $T=0$ einer Stufenfunktion an.
    \begin{equation*}
        f(E,T)=\frac{1}{e^{\frac{E-E_F}{k_BT}} + 1}
    \end{equation*}
    mit $E_F$: Fermienergie. Für $T=0K$ gilt:
    \begin{equation*}
        n=\int_0^{E_F}\left(\frac{2m}{\hbar^2}\right)^{\frac{3}{2}}\frac{\sqrt{E}}{2\pi^2}\diff E \Longrightarrow E_F=\frac{\hbar^2}{2m}(3\pi^2n)^{\frac32}
    \end{equation*}

    \textbf{Fermigeschwindigkeit} $v_F=\sqrt{\dfrac{2E_F}{m}}$
\end{sectionbox}


\subsection{Thermische Anregung}
Im quantenmechanischen werden elektronische Prozesse (z.B. Stromtransport, thermisches Aufheizen) als energetische Anregung von Elektronen aufgefasst.
Hierbei wird das Elektron auf ein höheres Energieniveau gebracht. 
Es können jedoch nur Elektronen nahe $E_F$ exitiert werden. 
Für $E<E_F$ gibt es keine freien Zustände und für $E>E_F$ gibt es keine Elektronen mehr. 
$n_{eff}$ beschreibt die Elektronendichte, die an thermischen Exitierungsprozessen teilnehmen können
\begin{equation*}
    n_{eff} \approx 2D(E_F)k_BT \approx \frac{3n}{2E_F}
\end{equation*}
\begin{equation*}
    C_{el}\approx\frac{6nk_B^2T}{E_F} \quad v_F=\sqrt{\frac{2E_F}{m}}
\end{equation*}

\subsection{Elektronen im periodischen Potential $U(x+a)=U(x)$}
Lösung der Schrödinger-Gleichung in der Form $\Psi_k(x)=u_k(x)e^{ikx}$ mit einer gitterperiodischen Funktion (Gitterkonstante $a$)
\begin{description}
    \item[erlaubte, diskrete Werte für die Block-Wellenzahl] $k=\dfrac{2\pi n}{Na}$
    \item[Wellen-Energie Relation] $\cos(ka) = \cos(ga) + \frac{m\beta a}{\hbar^2}\frac{\sin(ga)}{ga}, \; g=\frac{\sqrt{2mE}}{\hbar} \Rightarrow$ Nur bestimmte Bereiche für Energien erlaubt $\Rightarrow$ Energiebänder
    \item[Brillouinzone] Erste Zone im Bandstrukturdiagramm mit $k\in[-\frac{\pi}{2}, \frac{\pi}{2}]$ \\
    Gruppengeschwindigkeit wird an den Grenzen der Brillouinzone mit $k=\pm\frac{n\pi}{a}$ Null $\Rightarrow$ Energielücken
\end{description}